% ============================================================================
% Supplemental Material: GFD Letter (derivation, uniqueness, robustness)
% Accompanies: Gravity Field Dynamics: Decoding mass from observed rotation curves
% ============================================================================
\documentclass[aps,amsmath,amssymb,reprint,nofootinbib]{revtex4-2}

\usepackage[hidelinks]{hyperref}

% Single clean rule above references (same as letter.tex); refs in column 1 only on page 2
\makeatletter
\renewcommand{\bibsection}{%
  \par\vspace{1ex}%
  \noindent\rule{\linewidth}{0.4pt}%
  \par\vspace{0.5em}%
}
\let\@oldthebibliography\thebibliography
\let\@oldendthebibliography\endthebibliography
\renewcommand{\thebibliography}[1]{\begin{minipage}[t]{\columnwidth}\@oldthebibliography{#1}}
\renewcommand{\endthebibliography}{\@oldendthebibliography\end{minipage}}
\makeatother

\begin{document}

\title{Supplemental Material: Gravity Field Dynamics: Decoding mass from observed rotation curves}

\author{Stephen Nelson}
\email{stephen@fielddynamics.org}
\affiliation{Field Dynamics Research, ORCID: 0009-0008-6190-3269}

\date{February 2026}

\maketitle

\noindent\textbf{Capsule Summary.}
This document establishes the uniqueness of the constitutive law under field-closure postulates, states the relation of the Letter equations to the parent covariant theory, and provides sensitivity analysis for the kinematic inversion.

\section{Uniqueness of the constitutive law}
\label{sec:uniqueness}

The main Letter uses the relation $\mu(x) = x/(1+x)$ with $x = g_{\mathrm{t}}/a_0$, which in form coincides with the ``simple'' MOND interpolating function used in many rotation-curve fits. The important distinction is one of status, not shape. In phenomenological MOND, $\mu$ is typically chosen from a family of interpolating functions. Here it is not an ansatz: it is uniquely fixed by field-closure postulates at the boundary $\Sigma$.

Three postulates determine the constitutive response: (A1) local closure at the field boundary (equations close without external degrees of freedom); (A2) lossless transmission (no extra dissipation channel); (A3) a single calibration scale ($a_0$). Under these postulates, $\mu(x) = x/(1+x)$ is the unique solution. Other forms (e.g.\ $\mu = x/\sqrt{1+x^2}$, the ``standard'' MOND interpolant) correspond to different closure assumptions and remain testable against data. So $\mu$ is not a fitted choice; it is the unique form under these postulates.

A second route to the same result uses detailed balance and symmetry at the closure boundary. The field partitions into a resolved fraction $\varepsilon$ (transmitted) and an unresolved fraction $1 - \varepsilon$ (saturated). At equilibrium, the resolution rate scales as $x(1-\varepsilon)$ and the relaxation rate as $\varepsilon$; balancing them gives $\varepsilon = x/(1+x)$, so $\mu(x) \equiv \varepsilon(x) = x/(1+x)$.

It remains to fix the constant in the denominator. In the ghost-free bimetric theory, the two metric sectors enter the interaction potential symmetrically \cite{hassan2012}. The interaction vertex --- where both metrics couple at a single spacetime event --- contributes an additive structural constant to the constitutive law. If the constitutive law respects this sector symmetry so that the resolved and unresolved fractions are complementary under exchange, $\mu(x) + \mu(1/x) = 1$ with $x = g_{\mathrm{t}}/a_0$, then for the one-parameter family $\mu(x) = x/(n+x)$, this holds if and only if $n = 1$.

Our constitutive law therefore coincides with the commonly used ``simple'' MOND interpolant, but here it is uniquely determined by field-closure postulates rather than chosen. The claim is: unique under postulates A1, A2, and A3. This result explains retrospectively why the ``simple'' interpolant performs well in rotation-curve fits.

\section{Relation to covariant theory}
\label{sec:limit}

The two equations used in the Letter (Poisson and algebraic transition) are the spherically symmetric limit of a covariant scalar-tensor theory constructed within the tetrad formulation of general relativity and ghost-free bimetric gravity (\cite{mtw1973,wald1984}, \cite{hassan2012}). In that theory, the kinetic function $F$ in the scalar-tensor action (Letter Eq.\ 2) is fixed by the constitutive law: $\mu(\sqrt{y}) = dF/dy$ with $y = |\nabla\Phi|^2/a_0^2$, yielding $F(y) = y - 2\sqrt{y} + 2\ln(1+\sqrt{y})$. No free parameters beyond $G$ and $a_0$ enter the spherically symmetric limit. The full covariant derivation, including the coupling structure between the two metric sectors and the determination of $F$ from the bimetric action, will be presented in a dedicated paper. In this Letter the content under evaluation is the two field equations and the inverse pipeline they define.

\section{Uncertainty and robustness}
\label{sec:robustness}

The inverse pipeline uses $M_{\mathrm{dec}}(r) = r v^4 / (G(a_0 r + v^2))$. The $v^4$ factor amplifies measurement noise in the observed velocity. To reduce point-to-point scatter without changing the large-scale profile, a single fixed Gaussian smoothing (width equal to 6.2\% of the radial span) is applied to the observed velocities before decoding. The value 6.2\% is fixed identically for every galaxy and is not tuned per object.

How sensitive are the decoded mass profiles to this choice? Pipeline runs with smoothing fraction 4\%, 6.2\%, and 10\% give: for the Milky Way at 25 kpc, $M_{\mathrm{dec}}$ in the range $\sim 6.34\times10^{10}$ to $\sim 6.95\times10^{10}\,M_\odot$ (range about $+9.6\%$); at 50 kpc, $\sim 5.80\times10^{10}$ to $\sim 6.10\times10^{10}\,M_\odot$ (range about $+5.2\%$). For M33 at 17 kpc, $M_{\mathrm{dec}}$ runs from $\sim 1.35\times10^{10}$ to $\sim 1.39\times10^{10}\,M_\odot$ (range about $-3.1\%$). Varying the smoothing fraction from 4\% to 10\% therefore shifts the decoded masses by $\le 10\%$ (MW) and $\le 3\%$ (M33) at the reference radii. The decoded mass profiles and their comparison to baryonic mass are stable over this range.

At the outermost measured radii, where velocity uncertainties are largest, the decoded effective mass remains $6.4\sigma$ (MW, 50 kpc) and $17.2\sigma$ (M33, 17 kpc) below the $\Lambda$CDM enclosed mass. The main conclusion (effective mass consistent with baryonic mass and far below the $\Lambda$CDM mass budget) is thus robust to the smoothing choice and is not an artifact of the 6.2\% value.

\bibliography{../letter}

\end{document}
